\documentclass[12pt,a4paper]{article}
\usepackage{fontspec}
\setmainfont{Times New Roman}
\usepackage[top=2cm,bottom=2cm,left=2.5cm,right=2cm]{geometry}
\usepackage{enumitem}
\usepackage{tabularx}
\usepackage{booktabs}
\usepackage{array}
\newcolumntype{Y}{>{\centering\arraybackslash}X}
\pagestyle{plain}
\linespread{1.15}

\begin{document}
% --- HEADER 2 CỘT CHUẨN NGHỊ ĐỊNH 30/2020/NĐ-CP ---
\noindent
\begin{minipage}[t]{0.45\textwidth}
    \begin{center}
        \fontsize{11pt}{13pt}\selectfont
        CÔNG AN THÀNH PHỐ HÀ NỘI\\
        \textbf{PHÒNG KỸ THUẬT HÌNH SỰ (PC09)}\\[-0.2em]
        \rule{3.2cm}{0.6pt}\\[0.25cm]
        \fontsize{11pt}{13pt}\selectfont Số: \textbf{218/BC-PC09}
    \end{center}
\end{minipage}
\hfill
\begin{minipage}[t]{0.52\textwidth}
    \begin{center}
        \fontsize{11pt}{13pt}\selectfont
        \textbf{CỘNG HÒA XÃ HỘI CHỦ NGHĨA VIỆT NAM}\\
        \textbf{Độc lập - Tự do - Hạnh phúc}\\[-0.2em]
        \rule{3.6cm}{0.6pt}\\[0.25cm]
        \textit{Hà Nội, ngày 25 tháng 07 năm 2016}
    \end{center}
\end{minipage}

\vspace{0.6cm}

\begin{center}
    \textbf{\Large BÁO CÁO KẾT QUẢ KHÁM NGHIỆM TỬ THI SƠ BỘ}\\[0.15cm]
    \textit{}
\end{center}

\vspace{0.4cm}

\textbf{Kính gửi:}

\begin{itemize}[leftmargin=1.2cm, itemsep=2pt]
  \item Ban Giám đốc Công an Thành phố Hà Nội;
  \item Thủ trưởng Cơ quan Cảnh sát Điều tra – Công an TP. Hà Nội;
  \item Phòng Cảnh sát Hình sự (PC02) – Công an TP. Hà Nội.

\end{itemize}
Hồi 07 giờ 30 phút ngày 25 tháng 07 năm 2016, Phòng Kỹ thuật Hình sự (PC09) phối hợp với Cơ quan Điều tra tiến hành khám nghiệm tử thi nạn nhân trong vụ việc phát hiện chết người tại số 14 Đường Bờ Sông, Phường Phân khu Cảng, TP. Hà Nội.




\subsection*{I. LÝ LỊCH NẠN NHÂN}

\begin{itemize}[leftmargin=1.2cm, itemsep=2pt]
  \item \textbf{Họ và tên:} \textbf{NGUYỄN VĂN KHANG}
  \item \textbf{Sinh năm:} 1988
  \item \textbf{Địa chỉ thường trú:} Số 14 Đường Bờ Sông, Phường Phân khu Cảng, TP. Hà Nội.


\end{itemize}

\subsection*{II. TÌNH TRẠNG TỬ THI VÀ DẤU VẾT THƯƠNG TÍCH}


\begin{itemize}[leftmargin=1.2cm, itemsep=2pt]
  \item \textbf{Hiện tượng biến đổi tử thi:}
  \item Tử thi nam giới, thể trạng bình thường.
  \item Các khớp lớn nhỏ co cứng toàn diện; hoen ố tử thi tập trung cố định ở mặt sau lưng và mông, màu đỏ tím sẫm, ấn ngón tay không biến mất; giác mạc đục nhẹ.
  \item \textbf{Ước tính thời gian tử vong:} Căn cứ vào diễn biến của hoen tử thi và mức độ co cứng cơ bản, sơ bộ xác định nạn nhân tử vong vào khoảng thời gian từ \textbf{20 giờ 45 phút đến 21 giờ 15 phút đêm ngày 24/07/2016}.

  \item \textbf{Dấu vết thương tích bên ngoài:}
  \item \textbf{Vùng đầu - gáy:} Tại vùng chẩm gáy bên phải phát hiện một vết rách dập da kèm bầm tụ máu, kích thước đo được là \textbf{6,0 cm × 4,2 cm}, bờ mép vết rách dập không đều, đáy bầm tím. Cơ chế thương tích phù hợp với dạng tác động va đập cơ học mạnh vào vật tày có cạnh gờ \textit{(tương ứng với cạnh bàn gỗ tại hiện trường)}.
  \item \textbf{Vùng cổ (Vết thương chí mạng):} Tại vùng cổ bên trái phát hiện một vết thương hở dài, bờ mép vết thương sắc gọn, góc vết thương nhọn ăn sâu làm rách đứt hoàn toàn bó mạch cảnh \textit{(động mạch cảnh trái bị đứt mở rộng)}, gây xuất huyết ngoại nghiêm trọng.
  \item \textbf{Dấu hiệu bất thường trên da đầu:} Kiểm tra vùng thái dương và tóc mai bên trái ghi nhận: Một mảng tóc mai bên trái bị cắt sát da đầu, vết cắt bằng dụng cụ sắc nhọn, để lại mảng khuyết tóc rõ rệt trên diện tích khoảng \textbf{3cm × 2cm}.


\end{itemize}

\subsection*{III. NGUYÊN NHÂN TỬ VONG VÀ CƠ CHẾ HUNG KHÍ}


\begin{itemize}[leftmargin=1.2cm, itemsep=2pt]
  \item \textbf{Nguyên nhân chính dẫn đến tử vong:}
  \item Nạn nhân tử vong do \textbf{sốc mất máu cấp} hậu quả của vết thương hở làm rách đứt động mạch cảnh bên trái.
  \item Vết dập chấn thương vùng chẩm gáy phải \textit{(6.0cm × 4.2cm)} do va đập là tổn thương đi kèm, gây choáng hoặc mất khả năng tự vệ tức thời.

  \item \textbf{Nhận định hung khí gây thương tích:}
  \item Cơ chế vết rách đứt mạch máu vùng cổ trái phù hợp với tác động cứa/đâm bởi vật sắc nhọn, bằng chất liệu cứng.
  \item Khớp nối thương tích trực tiếp với vật chứng thu giữ tại hiện trường: Mảnh vỡ bình trà bằng gốm dính máu \textit{(ký hiệu 39)}.


\end{itemize}

\subsection*{IV. KIẾN NGHỊ NGHIỆP VỤ ĐIỀU TRA}


\begin{itemize}[leftmargin=1.2cm, itemsep=2pt]
  \item Chuyển giao ngay mẫu máu tử thi và mảnh vỡ bình trà \textit{(39)} để giám định gen (ADN) nhằm xác thực cơ chế hình thành dấu vết máu.


\end{itemize}

\vspace{0.8cm}
\noindent
\begin{tabularx}{\textwidth}{YY}
\toprule
\textbf{GIÁM ĐỊNH VIÊN / BÁC SĨ PHÁP Y} & \textbf{TRƯỞNG PHÒNG KỸ THUẬT HÌNH SỰ} \\
\midrule
\textit{(Ký, ghi rõ họ tên)} & \textit{(Ký tên, đóng dấu)} \\[1.5cm]
\textbf{Trần Tuấn Anh} & \textbf{Thượng tá Nguyễn Mạnh Hùng} \\
\bottomrule
\end{tabularx}
\end{document}
